\section{Sim-to-Real Predictivity Analysis}
\label{sec:predictivity}

This section fixes the analysis before the data exist, so that the reported
outcome cannot be a consequence of metric selection. Two questions are kept
separate throughout: how close simulated values are to physical ones, and
whether simulation orders methods the way water does. They can dissociate ---
a simulator can be biased and still rank correctly, or unbiased on average
and still invert an ordering --- and only the second bears on the simulator's
use as an evaluation instrument.

\subsection{Matched Conditions and Outcome Metrics}

A \emph{condition} is a tuple of avoidance method, obstacle, approach
geometry, and calibration level, run in both domains with the same nominal
route and speed. For each condition we record: success, collision, minimum
clearance, maximum lateral deviation, path length, maneuver duration,
recovery error to the nominal line, and control effort. These are the same
quantities the simulation campaigns already produce
(Sec.~\ref{sec:design}), so the comparison needs no domain-specific metric.

\subsection{Absolute Predictivity}

For a metric $m$ under condition $i$ at calibration level $S_k$, absolute
predictivity is the discrepancy between the simulated and physical values,
\begin{equation}
  e_{m}(i, S_k) \;=\; \bigl| \bar{m}^{\mathrm{sim}}(i, S_k)
                     - \bar{m}^{\mathrm{real}}(i) \bigr| ,
  \label{eq:abserr}
\end{equation}
aggregated over conditions and reported per level, with a scale-normalized
variant where metrics have different units. Since each side is a mean over
repetitions, both carry sampling uncertainty, and we report the gap with a
bootstrap interval rather than as a point value; a gap smaller than the
combined run-to-run spread is reported as indistinguishable from zero rather
than as agreement. Bounded outcomes such as success are compared as
proportions with Wilson intervals. This per-level, per-metric view is in the
spirit of judging sensor-model fidelity by its downstream
effect~\cite{mahajan2024quantify}.

\subsection{Ranking Predictivity and Its Statistical Requirements}

Rank preservation is the Kadian et al.\ question~\cite{kadian2020simrealp},
and adopting their Sim-vs-Real Correlation Coefficient is the natural choice
--- but adopting it uncritically would be a mistake in our design, and we
state the reason explicitly rather than reporting a number that cannot carry
weight. Their study correlates outcomes across nine navigation policies. Our
comparison contains two planners. A rank correlation over two items is
degenerate: it takes one of two values and has no meaningful null
distribution, so an SRCC computed over the planner pair alone would be
theatre.

Three defensible options exist, and we pre-commit to the following
resolution. \emph{First}, enlarge the ranked set along an axis we already
have: the ranked units are method \emph{stacks}, that is
estimator--planner combinations, since the estimator variants are
independently meaningful and already characterized (Sec.~\ref{sec:results}).
This yields a set large enough for a rank statistic without inventing methods
for the sole purpose of populating one. \emph{Second}, report
\emph{pairwise direction agreement} as the primary ranking result: for every
pair of stacks and every metric, whether simulation and reality agree on the
sign of the difference, with an exact binomial interval over pairs. This
statistic is well defined for two methods, degrades gracefully with ties, and
answers the question a practitioner actually asks --- if the simulator says A
is safer than B, is it? \emph{Third}, report rank correlation (Spearman or
Kendall, chosen by the tie structure) \emph{only} when the ranked set reaches
a pre-declared minimum size, and report it with its exact permutation
$p$-value rather than an asymptotic one.

We explicitly reject two temptations. We will not expand the experiment
solely to raise the item count for a correlation coefficient, because a
method added only to be ranked contributes noise, not evidence. And we will
not select, after unblinding, the metric on which ranking happens to be
preserved: the analysis reports every metric in
Sec.~\ref{sec:predictivity}\,A for every level, and the primary ranking
metric --- minimum clearance --- is declared here.

\subsection{Per-Level Attribution}

The ladder's payoff is the difference between adjacent levels. For each
metric we report absolute predictivity and ranking agreement at \SZERO{},
\SONE{}, \STWO{}, and \STHREE{}, and attribute any change between $S_{k-1}$
and $S_k$ to the class of mismatch that rung calibrated --- an attribution
that is only licensed because every other class is held fixed
(Table~\ref{tab:calibration}). A flat or adverse trajectory across levels is
a reportable outcome~\cite{truong2022rethinki}, and the pre-registration of
that interpretation is what keeps it reportable.

\subsection{Threats to Validity}

Three threats are anticipated. \emph{Confounding by repetition count}: the
physical campaign will have fewer repetitions than the simulated one, so
physical means are noisier; the bootstrap intervals of
Eq.~\eqref{eq:abserr} carry that asymmetry explicitly instead of hiding it.
\emph{Ceiling effects}: if every stack succeeds in both domains, ranking is
undefined on success and the analysis rests on continuous metrics ---
which is why the scenario set deliberately includes conditions that stress
clearance. \emph{Instrument error}: the external ground truth has its own
error budget, and a sim-to-real gap smaller than that budget is not a
measurement; the ground-truth error report (Sec.~\ref{sec:design}\,F) is
therefore a precondition for this analysis, not an appendix to it.

\SREAL
